\section{Discussion}
\label{sec:discussion}

Our findings distinguish social approval of minority judgments, the moral
reasons exchanged around them, and readers' consideration of alternative
judgments. These distinctions inform how online moral discussions can be
evaluated and presented to readers.

\subsection{Implications}
\label{subsec:discussion-findings}


Our findings suggest that examining the reasoning developed around dissent 
can help characterize the contribution of minority participation.
In high-agreement discussions, minority comments receive lower scores but attract more
replies and deeper threads than majority comments. Challenging replies
also frequently introduce moral reasons absent from their parent comments.
Opposition to a judgment can therefore coexist with the articulation of
additional reasons around it. Social approval, conversational engagement,
and the addition of moral reasons capture distinct dimensions of
participation. Moving beyond comment popularity, future evaluations
could examine which reasons are developed, disputed, and taken up
across an exchange.

Exposure to dissent may also be meaningful when readers agree with the
majority. In our survey, participants in \textit{Balanced}
reported lower judgment certainty and greater consideration of why
someone might reach a different conclusion than those in \textit{Highly Scored}.
These observations motivate evaluating minority influence through
engagement with competing reasons, consistent with research linking
minority influence to divergent thinking \cite{nemeth1987divergent}.
A reader's judgment alone does not reveal which alternatives they
considered or how they weighed them. Evaluations should therefore
examine the reasons readers give for their own and opposing judgments,
alongside judgment direction and certainty. This would make
consideration of alternative reasons an explicit outcome without
treating opinion change as the sole criterion of success.

These findings motivate presenting dissent with its reply context.
More than half of participants preferred seeing one majority comment and 
one minority comment rather than two majority comments.
Our reply analysis provides a further rationale for retaining the surrounding
exchange: challenges can introduce reasons absent from the minority
comment itself. Showing comments with their replies would allow readers
to inspect how judgments are justified and contested. Thread selection
should therefore consider the moral reasons represented and their
overlap, since opposing judgments do not necessarily offer distinct
reasoning. Future experiments could test which combinations of comments
and replies make alternative reasons accessible, while assessing
readers' experiences and contributors' willingness to continue
participating.

\subsection{Limitations and Future Work}
\label{subsec:discussion-limitations}

The observational findings concern Reddit moral-judgment communities and
may reflect story composition, unobserved exposure, and unequal follow-up.
Minority status is assigned retrospectively, and temporal profiles compare
successive comments rather than within-person change. Reconstructing the
opinion environment visible at posting and matching follow-up windows would
help clarify these reception patterns.

LLM annotations and the categories learned for each post need further
independent human validation. Addition labels capture differences from the
parent, without establishing argument quality or novelty across the
discussion. Tracking reasons through nested replies could reveal their
uptake and subsequent development.

The exploratory survey includes $80$ participants and four scenarios,
varying content, replies, and scores together. Responses are self-reported,
baseline judgments were not measured, alignment groups are defined after
exposure, and tests use unadjusted $p$-values. Larger studies with
prespecified comparisons should vary content and social cues independently.
Open-ended and delayed measures could assess which reasons readers use
and retain beyond their reported reflection.